\noindent
\begin{tabular}{ll}
\textbf{Thesis} & Trident: Tactical Research, Inspection, and Debris ENgagement Tool (UUV) \\
\textbf{Student} & Varit Roongrotekarnkha, Nabhatara Tongsri, Yotin Limyotin, Navavee Jungwongsuk \\
\textbf{Student ID.} & 65011620, 65011370, 65011648, 65011403 \\
\textbf{Degree} & Bachelor of Engineering \\
\textbf{Program} & Robotics and AI Engineering \\
\textbf{Year} & 2026 \\
\textbf{Thesis Advisor} & Asst. Prof. Dr. Jutarut Chaoraingern \\
\textbf{Co-Thesis Advisor} & \\
\end{tabular}

\vspace{1cm}

\begin{center}
\textbf{\Large ABSTRACT IN ENGLISH}
\end{center}

\vspace{0cm}

Thailand’s marine ecosystems are increasingly threatened by seabed pollution,
habitat degradation, and insufficient infrastructure inspection due to the high cost and
inaccessibility of traditional underwater monitoring systems. This project presents TRIDENT, a
modular, low-cost Unmanned Underwater Vehicle (UUV) developed to support marine
research, debris removal, and infrastructure maintenance in coastal regions. TRIDENT
integrates a vision-based system for marine debris detection, a robotic arm module for
interaction tasks, and high-resolution imaging for biodiversity documentation. Built using moldcast epoxy resin and silicone for cost-effective waterproofing, the vehicle incorporates an
efficient ballast tank system and pressure sensors for quasi-static depth control. Realtime data
acquisition is enabled via LAN communication with a Raspberry Pi and Arduino-based control
architecture running ROS. With a compact and energy-efficient design, TRIDENT enhances
underwater exploration in confined environments while addressing II critical marine
challenges. The system contributes to sustainable ocean monitoring by reducing operational
costs, improving inspection accessibility, and enabling data-driven conservation efforts.